\documentclass[a4paper,titlepage]{article}

%PACKAGES
\usepackage[utf8]{inputenc}
\usepackage[T1]{fontenc}
\usepackage[english]{babel}
\usepackage{amsmath}
\usepackage{amsthm}
\usepackage{amssymb}
\usepackage{mathrsfs}
\usepackage{fancyhdr}
\usepackage{lmodern}
\usepackage{graphicx}
\usepackage{geometry}
\usepackage{fancybox}
\usepackage{textcomp}

%Symbole euro
\usepackage{eurosym}

%Listings : affichage code
\usepackage{listings}


%Elements de la page de garde
\begin{document}

\begin{titlepage}

\begin{figure}
\centering
\includegraphics[width=5cm]{logo-ulg.png}
\end{figure}



\title{
\vspace{0.2cm}
\LARGE{\textbf{Second part report}} \\ \textsc{Discrete optimization}
\author{\textbf{Fabrice Servais} \small{(S111093})\\\textbf{Laurent Vanosmael} \small{(S114351}) \\ \\ Masters in Computer Science and Engineering - $1^{\text{st}}$ year}
\date{December 7, 2014}
\rule{15cm}{1.5pt}
}

%\geometry{hmargin=2.5cm}
\end{titlepage}

%DOCUMENT
\pagestyle{fancy}
\lhead{Second part report}
\rhead{Discrete optimization}

%Page de garde
\maketitle

\section{Introduction}



\section{Data}

The data listed below are the ones from the Integer Model that we have use in the first report.
\begin{itemize}		
	\item $sl_{k} = \left\{ 
	\begin{array}{ll}
		1 & \text{if the key $k$ is on the left side,}\\
		0 & \text{otherwise.}\\
	\end{array}
	\right.$

	\item $sr_{k} = \left\{ 
	\begin{array}{ll}
		1 & \text{if the key $k$ is on the right side,}\\
		0 & \text{otherwise.}\\
	\end{array}
	\right.$

	\item $fr_k$ : probability of the apparition of the letter $l$ in the considered language.

	\item $v_{l} = \left\{ 
	\begin{array}{ll}
		1 & \text{if the letter $l$ is a vowel,}\\
		0 & \text{otherwise.}\\
	\end{array}
	\right.$

	\item $V$ : number of vowels in the language.

	\item $w_{i,j}$ : probability that the letter $j$ follows the letter $i$ in a word. 

	\item $ks_{k}$ : strength of the finger associated to key $k$ ($ks_k \in [0,1])$.

	\item $d_k$ : distance that the finger attributed to the key $k$ has to cross to reach that key.
\end{itemize}

Data's Collect : \\

Concerning the data related to the keyboard ($sl_{k}$, $sr_{k}$, $d_{k}$) we have based our analysis on the keyboard present on the BEPO graphic in the assignment paper. \\
The distance is calculated as the number of "moves" from the home row the finger has to do. There is no diagonal moves. \\

For the strength of the finger we have arbitrary decided that thumb will be 0,  the index will  be 0.75 the middle finger 1, the annular 0.5 and the auricular 0.1. \\

We have match each key with a finger, so with the strength of the finger associated we can construct $ks_{k}$. \\

To collect the frequencies of the letters apparition and the bi-grams we have use the NGramExtraction tool from Deception.\\
To obtain the alphabet that we will use we just use the 47 letters and characters that have the highest number of occurrence. (47 is the number of key in our keyboard)\\
Then we "norm" the number of occurrence of each letter by the total to have the frequency of each letter.\\
We also norm the bi-gram so that the sum of each bi-gram frequency starting by a specified letter make always 1.\\

Data file organisation : \\
The file of data (named datas) is structured to be human readable, the first word is the name of the variable then the content (note: it has to be in this order for the parser program).
The data do not have the same format as in the model, for example it's the number of the key that is indicated in the sl, sr vectors. The parser translate it into binary data. We have do this in that way to avoid errors and to keep the data file readable.
fs represent the finger strength and the kbX the key associated to a certain finger (no key are for the thumbs).






\section{Model}

	\subsection{Variables}

\begin{itemize}
 	\item$kb_{k,l} = \left\{ 
	\begin{array}{ll}
		1 & \text{if the letter $l$ is on the key $k$,}\\
		0 & \text{otherwise.}\\
	\end{array}
	\right.$

	\item $vl = \left\{ 
	\begin{array}{ll}
		1 & \text{if the vowels are on the left side,}\\
		0 & \text{otherwise.}\\
	\end{array}
	\right.$

	\item $a_{i,j} = \left\{ 
	\begin{array}{ll}
		1 & \text{if to type $i$ and then $j$, it is not the same hand that is used,}\\
		0 & \text{otherwise.}\\
	\end{array}
	\right.$	
\end{itemize}

	\subsection{Constraints}

\begin{itemize}
	\item $\sum_l fr_l \cdot \left(\sum_k kb_{k,l} \cdot (sr_k - sl_k) \right) \geq 0$

	\item $\sum_k kb_{k,l} = 1 \text{, } \forall l$

	\item $\sum_l kb_{k,l} = 1 \text{, } \forall k$

	\item $\sum_l v_l \cdot \left( \sum_k kb_{k,l} \cdot sl_k \right) = V \cdot vl$

	\item $a_{i,j} = a_{j,i} \text{, } \forall i,j$

	\item $a_{i,i} = 0 \text{, } \forall i$

	\item $a_{i,j} \leq \sum_k kb_{k,i} \cdot sl_k + \sum_k kb_{k,j} \cdot sl_k$

	\item $a_{i,j} \geq \sum_k kb_{k,i} \cdot sl_k - \sum_k kb_{k,j} \cdot sl_k$

	\item $a_{i,j} \geq \sum_k kb_{k,j} \cdot sl_k - \sum_k kb_{k,i} \cdot sl_k$

	\item $a_{i,j} \leq 2 - \sum_k kb_{k,i} \cdot sl_k - \sum_k kb_{k,j} \cdot sl_k$
\end{itemize}

	\subsection{Objective function}

\[
	\min \left[ \sum_l fr_l \cdot \left( \sum_k ks_k \cdot d_k \cdot kb_{k,l}  \right) + \sum_i \sum_j w_{i,j} \cdot \left(1 - a_{i,j}\right) + \sum_i \sum_j w_{i,j} \left( \sum_k d_k \cdot kb_{k,i} + \sum_l d_l \cdot kb_{k,j} \right) \right]
\]



\section{Method}
To obtain an optimal keyboard related to our model by using the \textbf{row generation} method. To implement that method, we use 2 concepts that are present in Gurobi : the lazy constraint and the callback.
\paragraph{}
The lazy constraints are the constraints that are less likely to be violated, so they are not contained in the set of the beginning but are added as we go along the execution of the solver. In order to do that, we use the callback system. We provide a class containing a \texttt{callback()} function which is called at key steps of the resolution.
In our implementation, the function is called when the solver found a solution. This function checks if one of the lazy constraints is violated or not. If yes, we add it to the model by using the \texttt{addLazy()} function.
\paragraph{}
The constraints that are considered lazy are these one : $\forall i,j$,
\begin{itemize}
	\item $a_{i,j} \leq \sum_k kb_{k,i} \cdot sl_k + \sum_k kb_{k,j} \cdot sl_k$

	\item $a_{i,j} \geq \sum_k kb_{k,i} \cdot sl_k - \sum_k kb_{k,j} \cdot sl_k$

	\item $a_{i,j} \geq \sum_k kb_{k,j} \cdot sl_k - \sum_k kb_{k,i} \cdot sl_k$

	\item $a_{i,j} \leq 2 - \sum_k kb_{k,i} \cdot sl_k - \sum_k kb_{k,j} \cdot sl_k$
\end{itemize}
\paragraph{}



\section{Results}


\section{Improvements}

\newpage

\end{document}








